\section{Results}
\label{sec:results}

\begin{table*}[t]
\centering
\caption{AgentDojo \AdBench{} replay: \AdBenign{} benign tasks and \AdAttacks{} hijacked-agent attacks (\AdGoals{} attack goals), all pooled over the four suites. Benign: share of tasks that run with no person (\emph{unattended}), share that complete when a person approves every request (\emph{with review}), and requests per task. Attack: share whose goal calls all run without a person (\emph{unreviewed}), reach a person, or are refused. Policies are generated from tool names and descriptions only.}
\label{tab:agentdojo}
\begin{tabular}{lcccccc}
\toprule
& \multicolumn{3}{c}{Benign tasks} & \multicolumn{3}{c}{Attack goals} \\
\cmidrule(lr){2-4}\cmidrule(lr){5-7}
Policy & Unattended & With review & Requests/task & Unreviewed & To a person & Refused \\
\midrule
Allow everything (no filter) & \AdAllowUnattended & \AdAllowCompletes & \AdAllowAsks & \AdAllowGoalUnreviewed & \AdAllowGoalPerson & \AdAllowGoalDenied \\
Ask a person for everything & \AdAskUnattended & \AdAskCompletes & \AdAskAsks & \AdAskGoalUnreviewed & \AdAskGoalPerson & \AdAskGoalDenied \\
Tiers from descriptions & \AdTiersUnattended & \AdTiersCompletes & \AdTiersAsks & \AdTiersGoalUnreviewed & \AdTiersGoalPerson & \AdTiersGoalDenied \\
\quad + flow, hinted labels & \AdFlowUnattended & \AdFlowCompletes & \AdFlowAsks & \AdFlowGoalUnreviewed & \AdFlowGoalPerson & \AdFlowGoalDenied \\
\quad + flow, conservative labels & \AdConsUnattended & \AdConsCompletes & \AdConsAsks & \AdConsGoalUnreviewed & \AdConsGoalPerson & \AdConsGoalDenied \\
\bottomrule
\end{tabular}
\end{table*}


\subsection{RQ1: Containing a hijacked agent}
\label{sec:res-agentdojo}
Table~\ref{tab:agentdojo} pools the four suites. Without a filter every attack goal runs; asking a person for everything stops them all at \AdAskAsks{} requests per benign task and leaves no task unattended. Tiers derived from descriptions stop all but \AdTiersGoalUnreviewed{} \AdTiersGoalUnreviewedCI{} of attack goals from running without a person, refuse \AdTiersGoalDenied{} outright (deletions and removals), and cost \AdTiersAsks{} \AdTiersAsksCI{} requests per benign task, \AdAskReduction{} fewer than asking for everything; \AdTiersUnattended{} of benign tasks run unattended and \AdTiersCompletes{} complete when a person approves the requests. The goals that run unreviewed are one Slack goal: visiting an attacker's website, a single \AdTiersUnreviewedGoalTool{} call that the hint rules classify as a read.

Flow labels from the description hints change nothing: the hints mark tools sensitive only when descriptions mention secrets, credentials or personal records, which AgentDojo's descriptions (``returns the emails'') never do, so the Rule of Two never fires. Conservative labels close the remaining goal (\AdConsGoalUnreviewed{} unreviewed; with zero of \AdGoals{} goals, the rule-of-three bound is \AdRuleOfThree{}) at \AdConsAsks{} \AdConsAsksCI{} requests per benign task. Table~\ref{tab:agentdojo-suites} shows that the burden concentrates where reads are followed by sends: Slack, whose tasks read channels and web pages and then post.

\begin{table}[t]
\centering
\caption{Per AgentDojo suite: tasks, attack replays and goals, then benign unattended share, requests per benign task and unreviewed attack goals, for tiers alone and with conservative flow labels.}
\label{tab:agentdojo-suites}
\setlength{\tabcolsep}{3pt}
\resizebox{\columnwidth}{!}{%
\begin{tabular}{lrrrcccccc}
\toprule
& & & & \multicolumn{3}{c}{Tiers} & \multicolumn{3}{c}{+ flow (conservative)} \\
\cmidrule(lr){5-7}\cmidrule(lr){8-10}
Suite & Tasks & Attacks & Goals & Unatt. & Req. & Unrev. & Unatt. & Req. & Unrev. \\
\midrule
Workspace & 40 & 240 & 6 & 45\% & 0.65 & 0\% & 42\% & 0.68 & 0\% \\
Travel & 20 & 120 & 6 & 70\% & 0.30 & 0\% & 70\% & 0.30 & 0\% \\
Banking & 16 & 144 & 9 & 25\% & 0.88 & 0\% & 25\% & 0.88 & 0\% \\
Slack & 21 & 105 & 5 & 5\% & 1.62 & 20\% & 0\% & 2.48 & 0\% \\
\bottomrule
\end{tabular}}
\end{table}


\subsection{RQ2: Calibration is not containment}
\label{sec:res-deferral}
\begin{table}[t]
\centering
\caption{Calibrated deferral on the AgentDojo replay. \emph{Eligible}: share of benign and attack-goal approval requests that deferral could ever auto-approve. Right: benign tasks unattended and attack goals unreviewed if a calibrated threshold covered every eligible request, the case for a deployment whose history has few denials.}
\label{tab:deferral}
\setlength{\tabcolsep}{2.5pt}
\small
\begin{tabular}{lcccc}
\toprule
& \multicolumn{2}{c}{Eligible} & \multicolumn{2}{c}{All eligible approved} \\
\cmidrule(lr){2-3}\cmidrule(lr){4-5}
Policy & Benign & Attack & Unattended & Unreviewed \\
\midrule
Tiers & \DefTiersBenignEligible & \DefTiersGoalEligible & \DefTiersBenignUnattended & \DefTiersGoalUnreviewed \\
+ flow, hinted & \DefFlowBenignEligible & \DefFlowGoalEligible & \DefFlowBenignUnattended & \DefFlowGoalUnreviewed \\
+ flow, conservative & \DefConsBenignEligible & \DefConsGoalEligible & \DefConsBenignUnattended & \DefConsGoalUnreviewed \\
\bottomrule
\end{tabular}
\end{table}

Deferral reduces requests by auto-approving those a person almost always approves. Its guarantee holds in simulation (Table~\ref{tab:simulation}): the true denial rate under the chosen threshold exceeded $\alpha$ in at most \SimDecMaxViolation{} of runs with the decile grid, within $\delta = 10\%$. The simulation also exposed a power plateau in our first implementation, which tested every distinct score from the lowest upward: the first test then uses the smallest group in which zero denials could be significant, and the fixed sequence stops at its first failure, so the share of runs that found any threshold stayed near \SimEveryReadySixteenHundred{} however long the history. Testing a grid of score deciles, still fixed by the scores alone, raises it from \SimDecReadyHundred{} at 100 decisions to \SimDecReadySixteenHundred{} at 1{,}600, and the auto-approved share approaches the best possible \SimOptimal{}.

The guarantee concerns reviewers, not attackers. On the replay (Table~\ref{tab:deferral}), \DefTiersGoalEligible{} of attack-goal requests are eligible under tiers alone: they are ordinary sends and bookings. A deployment whose history contains few denials will calibrate a threshold that covers them, and if it covered every eligible request, \DefTiersGoalUnreviewed{} of attack goals would run without a person. Conservative flow labels make exposure-driven sends ineligible: only \DefConsGoalEligible{} of attack-goal requests remain eligible and \DefConsGoalUnreviewed{} of goals could run unreviewed, at the price of fewer benign tasks becoming unattended (\DefConsBenignUnattended{} versus \DefTiersBenignUnattended{}). Information flow is what makes deferral safe to switch on.

\begin{table}[t]
\centering
\caption{Learn then Test ($\alpha=0.05$, $\delta=0.1$), \SimRuns{} simulated histories per size: share of runs that find a threshold (\emph{ready}), share in which the true denial rate under it exceeds $\alpha$ (\emph{violated}), and the auto-approved share. The best threshold auto-approves \SimOptimal{}.}
\label{tab:simulation}
\setlength{\tabcolsep}{3pt}
\small
\begin{tabular}{rccccc}
\toprule
& \multicolumn{2}{c}{Every score} & \multicolumn{3}{c}{Decile grid} \\
\cmidrule(lr){2-3}\cmidrule(lr){4-6}
History & Ready & Violated & Ready & Violated & Auto share \\
\midrule
100 & 37\% & 3.4\% & 20\% & 0.9\% & 51\% \\
200 & 64\% & 3.8\% & 53\% & 0.1\% & 41\% \\
400 & 64\% & 2.7\% & 79\% & 0.2\% & 45\% \\
800 & 66\% & 3.2\% & 82\% & 0.1\% & 47\% \\
1600 & 63\% & 4.6\% & 98\% & 0.1\% & 49\% \\
\bottomrule
\end{tabular}
\end{table}


\subsection{RQ3: Typed actions, flow and prediction}
\begin{table}[t]
\centering
\caption{Built-in labelled suites (written alongside the packs; regression tests, not an independent benchmark): exact-verdict accuracy of the full pack, of tiers alone, and the single layer whose removal costs most.}
\label{tab:internal}
\setlength{\tabcolsep}{3pt}
\resizebox{\columnwidth}{!}{%
\begin{tabular}{lrccl}
\toprule
Suite & Cases & Full & Tiers only & Costliest removal \\
\midrule
Agent hosts (coding) & 22 & 100\% & 59\% & argument rules (68\%) \\
Authorised testing & 14 & 100\% & 64\% & scope (71\%) \\
Coding / MCP & 22 & 100\% & 64\% & argument rules (64\%) \\
Flow sessions (coding) & 19 & 100\% & 100\% & none \\
Flow sessions (SOC) & 27 & 100\% & 67\% & information flow (67\%) \\
SOC response actions & 51 & 100\% & 80\% & prediction (92\%) \\
SOC triage & 33 & 100\% & 85\% & clamps and redaction (94\%) \\
Typed actions (SOC) & 12 & 100\% & 50\% & typed actions (50\%) \\
Typed actions (testing) & 32 & 100\% & 53\% & typed actions (59\%) \\
\bottomrule
\end{tabular}}
\end{table}

Table~\ref{tab:internal} summarises the built-in suites; every pack reaches 100\% on its suites by construction, so the informative numbers are the ablations and the earlier versions. The authorised-testing pack before typed actions got \BaTypedTestBefore{} of the typed suite right. It refused in-scope requests written as ranges, wildcard octets, \texttt{host:port}, a trailing dot or a comma list (false stops \BaTypedTestFalseStopBefore{}). It sent a gopher URL to an internal Redis server to a reviewer, and it judged \texttt{http://evil.example\textbackslash{}@web.lab.example/} to be in scope because Python's URL parser reports the host after the \texttt{@}, while HTTP clients disagree about that host~\cite{orangetsai}. The SOC pack before typed actions got \BaTypedSocBefore{} right: blocking \texttt{0.0.0.0/1}, \texttt{10.0.0.0/8}, a 65{,}535-address range or \texttt{*.com}, and internal addresses written in hex or as an integer, all reached a person instead of being refused. On the multi-step SOC flow suite the pack before information flow got \BaFlowBefore{} right: it sent every message, internal hand-offs included, to a person (false stops \BaFlowFalseStopBefore{}, counting requests as stops) while letting a payroll attachment go to a public sandbox. On labelled SOC response actions, prediction against the inventory moves inappropriate containment from reaching a person (\PredInapPersonWithout{} without prediction) to being refused with a stated prediction (\PredInapPersonWith{} still reach a person), without refusing any additional appropriate action (\PredAppDeniedWith{} with and \PredAppDeniedWithout{} without).

\subsection{RQ4: Authoring from failures}
On the AgentDojo replay we label benign calls as acceptable when they are allowed, modified or asked, and attack-goal calls when they are asked or denied, then propose edits on one half of the tasks and goals and test on the other half (both directions, general strategy). On Slack, edits proposed from one half close the unreviewed goal in the other (\AuSlackGoalBefore{} $\to$ \AuSlackGoalAfter{}), but they do so by asking for reads that precede attack goals, and benign requests per task rise from \AuSlackAsksBefore{} to \AuSlackAsksAfter{}. Authoring optimises the labels it is given: call-level labels that only penalise unreviewed attack calls and refusals push edits toward asking. A burden-aware objective, or labels on the effectful step of an attack only, is needed before such edits are adopted; the tool reports the burden so that a person can see it.

\subsection{Overhead}
The decision itself takes a median of \LatMedian{}\,ms (95th percentile \LatPNinetyFive{}\,ms) per call in process, excluding network transport. The full replay, including all policies, deferral analysis and authoring folds, runs in \AdSeconds{}\,s in a single Python process.
